
\section{\BookRecordIndex{T0058}半正性}\label{sec:6-4}
\zhindex{banzhengxing}{半正性}

本节考察一些能够控制 $\overline{\mathcal M}_{0,k}(A;J)$ 较高\BookRecordIndex{T0072}层维数的假设。这些条件的
关键，是能够由 $J$-全纯球面表示的同调类 $A\in H_2(M;\mathbb Z)$ 的第一陈数
\[
  c_1(A)=\langle c_1(TM,J),A\rangle
\]
所取的值。事实上，如本章导言所述，我们需要一些条件，以保证对一般的 $J$ 不存在陈数为
负的 $J$-全纯球面。

若存在数 $\tau>0$，使得对每个 $A\in\pi_2(M)$ 都有
\[
  c_1(A)=\tau\omega(A),
\]
就称\BookRecordIndex{T0358}辛流形 $(M,\omega)$ 是\emph{（球面）单调的}；若上述等式对某个 $\tau\ge0$ 成立，
则称它是\emph{\BookRecordIndex{T0056}半单调的}。这是一个自然但非常强的条件。例如，考虑 \BookRecordIndex{T0017}Fano 簇，即反\BookRecordIndex{T0103}典范丛
$K^*=\Lambda^{n,0}TM$ 丰富的复流形 $(M,J)$。在单连通情形，这等价于存在一个 凯勒
度量，使其 凯勒 形式 $\omega$ 是\BookRecordIndex{T0284}球面单调的。
\footnote{一般而言，\BookRecordIndex{T0207}凯勒 流形 $(M,J)$ 的反\BookRecordIndex{T0103}典范丛丰富，当且仅当 $c_1(TM)$ 类中存在
凯勒 形式。因此，对所有同调类 $A$，而不仅是\BookRecordIndex{T0285}球面类，都有
$[\omega](A)=c_1(A)$。}
\sourcefootnote{原书脚注 1 把“反\BookRecordIndex{T0103}典范丛”误写成“\BookRecordIndex{T0103}典范丛”；但正文刚定义
$K^*=\Lambda^{n,0}TM$，且脚注随后要求 凯勒 类等于 $c_1(TM)$，两者都唯一确定此处
应为反\BookRecordIndex{T0103}典范丛。}
对我们而言，更切合需要的是下面这个专为当前用途而设、但很有用的定义。

\begin{definition}\label{def:6-4-1}
若对每个 $A\in\pi_2(M)$ 都有
\[
  \omega(A)>0,\qquad c_1(A)\ge3-n
  \qquad\Longrightarrow\qquad c_1(A)\ge0,
\]
就称\BookRecordIndex{T0358}辛流形 $(M,\omega)$ 是\emph{\BookRecordIndex{T0058}半正的}。
\end{definition}

因此，\BookRecordIndex{T0056}半单调\BookRecordIndex{T0358}辛流形是\BookRecordIndex{T0058}半正的；光滑的 NEF 射影簇（即数值有效簇）也是如此。
\footnote{本书第一版用“弱单调”一词代替“\BookRecordIndex{T0058}半正”。不过，现在的文献似乎更多使用后一个
名称。还要注意，因为这个定义只是为方便而采用，有时会依具体情形作相应修改。}
反过来，许多\BookRecordIndex{T0058}半正流形 $(M,\omega)$ 对某个\BookRecordIndex{T0363}辛形式 $\omega'$ 是\BookRecordIndex{T0056}半单调的，但 $\omega'$ 很
可能不同于 $\omega$。为了刻画\BookRecordIndex{T0058}半正流形，我们作如下定义。

\begin{definition}\label{def:6-4-2}
设 $(M,\omega)$ 是紧\BookRecordIndex{T0358}辛流形。$(M,\omega)$ 的\emph{\BookRecordIndex{T0418}最小陈数}是整数
\[
  N=\inf\bigl\{k>0\bigm|\text{存在 }A\in\pi_2(M),\ c_1(A)=k\bigr\}.
\]
若对每个 $A\in\pi_2(M)$ 都有 $c_1(A)=0$，则约定\BookRecordIndex{T0418}最小陈数为 $N=\infty$。若
$N\ne\infty$，则
\[
  \langle c_1,\pi_2(M)\rangle=N\mathbb Z.
\]
\end{definition}
\zhindex{zuixiaochenshu}{最小陈数}

\begin{exercise}\label{ex:6-4-3}
证明：\BookRecordIndex{T0358}辛流形 $(M,\omega)$ \BookRecordIndex{T0058}半正，当且仅当下列条件之一成立。
\begin{enumerate}[label=(\roman*),leftmargin=3em]
\item $(M,\omega)$ \BookRecordIndex{T0056}半单调；
\item 对每个 $A\in\pi_2(M)$，都有 $c_1(A)=0$；
\item \BookRecordIndex{T0418}最小陈数 $N$ 满足 $N\ge n-2$。
\end{enumerate}
这些条件并不互斥。我们特别强调 $c_1$ 在 $\pi_2(M)$ 上为零的情形，是因为它包含重要的
Calabi--Yau 流形情形；不过，这恰是 (iii) 中 $N=\infty$ 的情形。还要注意，(iii) 蕴含
每个维数不超过六的\BookRecordIndex{T0358}辛流形都是\BookRecordIndex{T0058}半正的。它还包含 $c_1$ 在 $\pi_2(M)$ 上的限制是
$[\omega]$ 的负倍数的某些情形。
\end{exercise}

下面的引理突出了\BookRecordIndex{T0058}半正性最重要的推论。

\begin{lemma}\label{lem:6-4-4}
设 $(M,\omega)$ 是\BookRecordIndex{T0058}半正\BookRecordIndex{T0358}辛流形，$J$ 是被 $\omega$ \BookRecordIndex{T0372}驯服的\BookRecordIndex{T0088}殆复结构，而
$u:S^2\to M$ 是 $J$-全纯球面。若 $D_u$ 的余核维数为零或一，则
$c_1(A)\ge0$，其中 $A\in\pi_2(M)$ 是 $u$ 的同伦类。
\end{lemma}

\begin{proof}
若 $u$ 是常值映射，则 $c_1(A)=0$，故设 $u$ 非常值。于是
$\omega(A)=E(u)>0$。此外，若 $\zeta\in\operatorname{Vect}(S^2)$ 是无穷小 \BookRecordIndex{T0040}Möbius
变换，则 $du\circ\zeta\in\ker D_u$，所以 $\dim\ker D_u\ge6$。由于
$\dim\operatorname{coker}D_u\le1$，这给出
\[
  \operatorname{index}D_u=2n+2c_1(A)\ge5.
\]
因而 $c_1(A)\ge3-n$；再由\BookRecordIndex{T0058}半正性假设，$c_1(A)\ge0$。
\end{proof}

条件 $(M,\omega)$ \BookRecordIndex{T0058}半正只依赖于上同调类 $[\omega]$ 和
$c_1=c_1(TM,J)$；但实际需要的是前一引理的结论，而它涉及具体\BookRecordIndex{T0088}殆复结构 $J$ 的性质。
因此，下面的定义很有用。

\begin{definition}\label{def:6-4-5}
设 $(M,\omega)$ 是\BookRecordIndex{T0358}辛流形，$\kappa>0$。若每个能量 $E(u)\le\kappa$ 的
$J$-全纯球面 $u:\mathbb{CP}^1\to M$ 都有非负陈数，就称被 $\omega$ \BookRecordIndex{T0372}驯服的\BookRecordIndex{T0088}殆复结构
$J\in\mathcal J_\tau(M,\omega)$ 是\emph{$\kappa$-\BookRecordIndex{T0058}半正的}。若 $J$ 对每个 $\kappa$
都是 $\kappa$-\BookRecordIndex{T0058}半正的，就称 $J$ 是\emph{\BookRecordIndex{T0058}半正的}；这就是说，每个 $J$-全纯球面都有
非负陈数。
\end{definition}

用 $\mathcal J_+(M,\omega;\kappa)$ 表示 $M$ 上所有被 $\omega$ \BookRecordIndex{T0372}驯服且
$\kappa$-\BookRecordIndex{T0058}半正的\BookRecordIndex{T0088}殆复结构之集，并记
\[
  \mathcal J_+(M,\omega)
  :=\bigcap_{\kappa>0}\mathcal J_+(M,\omega;\kappa)
\]
为所有被 $\omega$ \BookRecordIndex{T0372}驯服的\BookRecordIndex{T0058}半正\BookRecordIndex{T0088}殆复结构之集。（注意，这些集合可能为空。）

\begin{remark}\label{rem:6-4-6}
这里我们使用所有被 $\omega$ \BookRecordIndex{T0372}驯服的\BookRecordIndex{T0088}殆复结构组成的空间
$\mathcal J_\tau(M,\omega)$，所以上述定义也以这类 $J$ 表述。不过，若改为考虑所有与
$\omega$ \BookRecordIndex{T0348}相容的\BookRecordIndex{T0088}殆复结构组成的空间 $\mathcal J(M,\omega)$，并把各处“被 $\omega$
\BookRecordIndex{T0372}驯服”换成“与 $\omega$ \BookRecordIndex{T0348}相容”，本节的全部结果仍然成立。使用被 $\omega$ \BookRecordIndex{T0372}驯服的\BookRecordIndex{T0088}殆复
结构的好处是，可以在不改变 $J$ 的情况下改变 $\omega$。
\end{remark}

\begin{lemma}\label{lem:6-4-7}
设 $(M,\omega)$ 是紧辛 $2n$-流形。则对每个 $\kappa>0$，集合
$\mathcal J_+(M,\omega;\kappa)$ 在 $\mathcal J_\tau(M,\omega)$ 的
$C^\infty$ 拓扑下是开的。
\end{lemma}

\begin{proof}
设 $J_\nu\in\mathcal J_\tau(M,\omega)$ 是一列并非 $\kappa$-\BookRecordIndex{T0058}半正的\BookRecordIndex{T0088}殆复结构，并在
$C^\infty$ 拓扑下收敛到 $J\in\mathcal J_\tau(M,\omega)$。于是存在一列
$J_\nu$-全纯球面 $u_\nu:\mathbb{CP}^1\to M$，满足
\[
  c_1(u_\nu)<0,\qquad E(u_\nu)\le\kappa.
\]
由定理~\ref{thm:5-3-1}，$u_\nu$ 有一个子列\BookRecordIndex{T0141}格罗莫夫收敛到\BookRecordIndex{T0340}稳定映射
$\mathbf u$。由定理~\ref{thm:5-2-2}，它的某个分支 $u_\alpha$ 必有负陈数，而所有分支
都满足 $E(u_\alpha)\le\kappa$。因此 $J$ 不是 $\kappa$-\BookRecordIndex{T0058}半正的。
\end{proof}

即使 $J$ \BookRecordIndex{T0058}半正，由 $\kappa$-\BookRecordIndex{T0058}半正结构组成的 $J$ 的邻域大小仍可能依赖于 $\kappa$。
所以，紧\BookRecordIndex{T0358}辛流形上的\BookRecordIndex{T0058}半正\BookRecordIndex{T0088}殆复结构之集未必是开的；不过，它是剩余的。此外，它本身虽未必
道路连通，逼近它的集合 $\mathcal J_+(M,\omega;\kappa)$ 却道路连通，而这已经足够用于
我们的目的；见定理~\ref{thm:6-6-1} 的证明。

\begin{lemma}\label{lem:6-4-8}
设 $(M,\omega)$ 是紧\BookRecordIndex{T0058}半正\BookRecordIndex{T0358}辛流形。则 $\mathcal J_+(M,\omega)$ 在
$\mathcal J_\tau(M,\omega)$ 中稠密；而且对每个 $\kappa>0$，集合
$\mathcal J_+(M,\omega;\kappa)$ 都是开且道路连通的。
\end{lemma}

\begin{proof}
考虑所有满足下述条件的被 $\omega$ \BookRecordIndex{T0372}驯服的\BookRecordIndex{T0088}殆复结构 $J$ 组成的集合
$\mathcal J_r$：对每个 $J$-全纯球面 $u$，算子 $D_u$ 的余核维数至多为一。由
引理~\ref{lem:6-4-4}，$\mathcal J_r\subset\mathcal J_+(M,\omega)$；由
定理~\ref{thm:3-1-8}，$\mathcal J_r$ 是\BookRecordIndex{T0300}剩余集。因此
$\mathcal J_+(M,\omega)$ 在 $\mathcal J_\tau(M,\omega)$ 中稠密。

现在取 $\kappa>0$ 和
$J_0,J_1\in\mathcal J_+(M,\omega;\kappa)$。由于
$\mathcal J_+(M,\omega;\kappa)$ 在 $\mathcal J_\tau(M,\omega)$ 中是开的，可在其中
任意小地扰动端点，故不失一般性可设
$J_0,J_1\in\mathcal J_{\mathrm{reg}}(M,\omega)$。定理~\ref{thm:3-1-8} 断言，
$J_0$ 与 $J_1$ 可由定义~\ref{def:3-1-7} 意义下的\BookRecordIndex{T0407}正则同伦
$\{J_\lambda\}$ 连接。由\BookRecordIndex{T0407}正则同伦的定义，对每个 $\lambda$ 和每个
$J_\lambda$-全纯球面 $u:S^2\to M$，都有
$\dim\operatorname{coker}D_u\le1$。所以，对每个 $\lambda\in[0,1]$，
\[
  J_\lambda\in\mathcal J_r\subset\mathcal J_+(M,\omega).
\]
这证明了引理。
\end{proof}

最后要指出，有些有趣的\BookRecordIndex{T0358}辛流形并不\BookRecordIndex{T0058}半正，但对某些（并非所有）正则且被 $\omega$ \BookRecordIndex{T0372}驯服的
\BookRecordIndex{T0088}殆复结构，仍然有引理~\ref{lem:6-4-4} 的结论。在这种情形，本书发展的技术可以定义流形
$(M,\omega)$ 的不变量，而先验看来它可能依赖于 $J$ 的选择；见
注~\ref{rem:6-6-2}。若使用\BookRecordIndex{T0143}格罗莫夫--威滕不变量可对所有\BookRecordIndex{T0358}辛流形定义这一事实，就能\BookRecordIndex{T0326}推出
这些不变量与 $J$ 无关。但即使只在本书的框架内，我们仍能得到三元组 $(M,\omega,J)$ 的
有趣不变量。这在 凯勒 几何中特别重要，因为人们关心的往往更多是底\BookRecordIndex{T0072}层复流形
$(M,J)$ 的性质，而不是某个具体 凯勒 形式 $\omega$ 的选择。

举例来说，考虑射影空间的乘积
\[
  M=\mathbb{CP}^{n_1}\times\cdots\times\mathbb{CP}^{n_k},
\]
配备\BookRecordIndex{T0363}辛形式
\[
  \omega_\lambda=\sum_{i=1}^k\lambda_i\omega_i,
\]
其中 $\omega_i$ 是 $\mathbb{CP}^{n_i}$ 上通常 凯勒 形式到 $M$ 的拉回，并归一化为在
每条复直线上的积分等于 $\pi$。那么，$M$ 单调当且仅当所有比值
$\lambda_i/(1+n_i)$ 相等。在大多数其他情形，流形 $(M,\omega_\lambda)$ 甚至不是\BookRecordIndex{T0058}半正的。

\begin{exercise}\label{ex:6-4-9}
求乘积流形
$(\mathbb{CP}^{n_1}\times\cdots\times\mathbb{CP}^{n_k},\omega_\lambda)$ \BookRecordIndex{T0058}半正的
充要条件。另一方面，证明：若
$J_{\mathrm{split}}=J_1\times\cdots\times J_k$ 是各因子上\BookRecordIndex{T0088}殆复结构 $J_i$ 的乘积，
则每个非常值 $J_{\mathrm{split}}$-全纯球面都有正陈数。
\end{exercise}

\begin{exercise}\label{ex:6-4-10}
考虑流形
\[
  M=S^2\times S^2
\]
和\BookRecordIndex{T0363}辛形式
\[
  \omega_\lambda:=\lambda\pi_1^*\sigma\oplus\pi_2^*\sigma.
\]
（见例~\ref{ex:3-3-6}。）这个流形上的标准复结构是\BookRecordIndex{T0058}半正的。当 $\lambda=1$ 时，
$(S^2\times S^2,\omega_\lambda)$ 单调，所以每个与 $\omega_\lambda$ \BookRecordIndex{T0348}相容的 $J$ 都是正的，
即每条 \BookRecordIndex{T0003}$J$-全纯曲线都有正陈数。

现在把 $S^2$ 等同于 $\mathbb C\cup\{\infty\}$，使得
\[
  \sigma=\frac{dx\wedge dy}{(1+x^2+y^2)^2}
\]
（见习题~\ref{ex:4-2-1}）。对 $m>0$，定义嵌入
\[
  u_m:S^2\longrightarrow S^2\times S^2,
  \qquad u_m(z):=(z,\overline z^{\,m}).
\]
证明：$u_m$ 关于 $\omega_\lambda$ 是辛嵌入，当且仅当 $\lambda>m^2$。\BookRecordIndex{T0326}推出：对每个
$\lambda>m^2$，都存在一个与 $\omega_\lambda$ \BookRecordIndex{T0348}相容的 $J$，使它容许一条陈数为
$2-2m$ 的 $J$-全纯球面。
\sourcefootnote{原书把该球面的陈数写成 $2-m$；但
$[u_m]=(1,-m)$，而 $c_1(T(S^2\times S^2))=(2,2)$，故其陈数唯一为 $2-2m$。
同句又误称“Exercise 3.3.6”，实际对应例~\ref{ex:3-3-6}，一并改正。}
（事实上，如例~\ref{ex:3-3-6} 所说明的，只要 $\lambda>m$，上述结论就成立；但用目前
这些很初等的方法很难看出这一点。）
\end{exercise}
